\seoli{유효숫자 자리 적용의 일관성, 평균값 계산에서의 유효숫자 계산 규칙 이해 및 적용,
그리고 유효숫자의 필요성 인식이라는 세 영역}을 종합하면,
%
과학영재 학생들은 유효숫자 규칙을 개념적으로는 상당 수준 이해하고 있으나,
이를 실제 탐구 과정에서 일관되게 적용하는 데에는 어려움을 겪고 있음을 확인할 수 있다.
%
즉, 지식 수준의 이해는 존재하지만 실제 측정값 기록, 계산 과정, 필요성 판단으로
자연스럽게 연결되지 않았으며, 특히 측정 상황과 계산 상황에서는 `형식적 절차 적용'에
머무르는 모습이 두드러졌다.
%
학생들은 정밀도와 불확실도의 개념을 단순한 자릿수 규칙으로 해석하는 경향을 보였고,
탐구 맥락에 따라 필요성 인식이 달라지는 불안정한 판단을 나타냈다.
%
이는 유효숫자 개념이 탐구적 사고의 일부로 내면화되기보다,
규칙 중심의 표면적 지식으로 작동하고 있음을 의미한다.
%
다시 말해, 지식 수준에서는 이해하지만 실제 탐구 과정에서는 일관된 판단으로 연결되지 못하고 있으며,
이는 탐구적 사고의 단절로 이어지고 있다.